\chapter{Conclusiones y Trabajo Futuro}
\label{05_00_00_conclusiones}

Esta tesis propuso el PSBPM-FD, un modelo bayesiano no paramétrico para series de tiempo funcionales que estima la distribución condicional completa de la curva siguiente dado el pasado del proceso. Las curvas se representan por sus scores de componentes principales funcionales, la ley de cada score dado el estado rezagado es una mezcla probit stick-breaking de regresiones normales, y la predictiva de los scores se transporta a la curva por la expansión funcional. Sobre esa predictiva se definieron tres predictores puntuales, tres construcciones de intervalo y los comportamientos predictivos, que identifican en cada origen los desenlaces posibles de la curva siguiente y la probabilidad de cada uno, y de ellos un pronóstico y su cota. El modelo se evaluó en tres escenarios simulados y en el tráfico de una autopista urbana, con el FAR y el Random Forest como referencias. Este capítulo reúne lo que ese estudio establece, lo que deja abierto y las direcciones en que conviene continuarlo.

\section{Conclusiones}
\label{05_01_00_conclusiones}

La primera pregunta del estudio era en qué condiciones la predicción puntual del PSBPM-FD supera a la de los modelos de referencia. La respuesta es que la ventaja sobre el FAR es grande donde la media condicional es no lineal y pequeña pero sostenida donde es lineal o casi lineal. En el escenario con cambio de régimen la esperanza reduce el error del FAR en torno a un quinto y lo supera en todas las ventanas de prueba, mientras que con heterocedasticidad condicional y con multimodalidad en la ley de los scores la reducción es de un par de puntos porcentuales que se repite en la gran mayoría de las ventanas (Cuadros~\ref{tab:res_bloqueA_ganancia} y~\ref{tab:res_bloqueA_ventanas}). En el tráfico, donde la ley de la curva siguiente es una mezcla de perfiles, la ventaja depende del resumen, ya que la mediana puntual supera al FAR bajo MAE con una y con dos covariables rezagadas y la esperanza no lo supera bajo RMSE, porque promedia perfiles que no ocurren (Sección~\ref{04_03_01_error_puntual}). Entre los predictores, la esperanza es la mejor bajo RMSE y la mediana puntual bajo MAE donde la predictiva es multimodal, como corresponde a la pérdida que cada uno minimiza, y el átomo más probable es el más débil. El Random Forest, con las mismas covariables, es el mejor predictor puntual cuando la dinámica tiene la estructura de umbral que sus cortes reproducen y cuando solo se dispone del rezago de un día en el tráfico, y queda detrás del PSBPM-FD en los demás casos.\\

La segunda pregunta era en qué condiciones los intervalos de credibilidad ofrecen un mejor desempeño, y es donde el modelo muestra su ventaja más clara. Las bandas puntuales reducen el puntaje de Winkler del FAR en los tres escenarios y en las dos especificaciones del tráfico, con una reducción de un quinto donde hay cambio de régimen y de un décimo en el tráfico, y lo superan en la mayoría de las ventanas de prueba en todos los casos (Cuadros~\ref{tab:res_bloqueB_ganancia} y~\ref{tab:apl_bloqueB_ganancia}). La razón es la misma en la simulación y en los datos reales, el ancho de la banda del PSBPM-FD sigue al estado del proceso, angosto en los tramos tranquilos y ancho en los episodios de mayor varianza, en el régimen volátil o en los días cuyo perfil la información no decide, mientras que la banda del FAR hereda una varianza residual fija y tiene el mismo ancho en todos los orígenes (Secciones~\ref{03_08_02_intervalos} y~\ref{04_03_02_intervalos}). Esa adaptación es la que una ley de innovación rígida no puede dar y la razón por la que la ventaja en los intervalos aparece incluso donde la media condicional es lineal y el error puntual empata. Entre las construcciones, la equi-colas es la referencia, la región de máxima densidad solo se separa de ella donde la predictiva tiene valles y ahí la mejora, y la banda simultánea es la única calibrada para cubrir la curva completa, al precio de un ancho que le hace perder el Winkler.\\

La tercera pregunta era en qué medida los comportamientos de la predictiva reproducen los regímenes del proceso y permiten pronosticar a partir de ellos. En el escenario con cambio de régimen los comportamientos los producen los scores que cambian de régimen y no la dimensión $M$ (Sección~\ref{03_08_03_comportamientos}), y en el tráfico corresponden a los dos tipos de día, el laborable y el no laborable, aparecen donde el calendario deja ese tipo sin decidir y la cota muestra que uno de ellos es la curva que ocurre (Secciones~\ref{04_03_03_comportamientos} y~\ref{04_03_04_pronostico}). Lo que la información cambia no es ese agrupamiento sino la probabilidad que recibe cada grupo, de modo que al agregar el rezago semanal la predictiva se vuelve unimodal en los días en que esa información decide el perfil y conserva varios comportamientos donde no lo decide. El modelo no impone la multimodalidad, la produce cuando lo que sabe no basta para elegir entre estados. Pronosticar con el comportamiento más probable es una apuesta que gana al FAR con holgura, empata con la esperanza bajo MAE y pierde bajo RMSE, porque las elecciones equivocadas pesan más que las acertadas (Cuadros~\ref{tab:res_pronostico} y~\ref{tab:apl_pronostico_101}). Su cota muestra que la predictiva contiene casi siempre la curva que ocurre, es decir, que el modelo captura la ley condicional y no solo su media, y que la pérdida del pronóstico se produce al elegir, y la probabilidad del comportamiento elegido indica cuándo conviene hacerlo, de manera que con una probabilidad alta se pronostica con el comportamiento y con una baja se reportan las opciones (Cuadros~\ref{tab:res_acierto_p1} y~\ref{tab:apl_acierto_p1}).\\

En conjunto, lo que distingue al PSBPM-FD de un modelo lineal no es el error puntual sino que su predictiva tiene la forma de la ley condicional, y esa forma es la que da los intervalos adaptativos y los comportamientos. Su alcance queda delimitado por las otras cinco series reales del Anexo~\ref{aneD_00_00_otras_series}, sin un estado oculto que separe perfiles, donde la predictiva es unimodal y el modelo queda cerca del FAR. La ventaja aparece en series cuya ley condicional tiene más de un comportamiento o cuya dispersión cambia con el estado, y es en esas series donde el modelo debe usarse.

\section{Limitaciones}
\label{05_02_00_limitaciones}

Las conclusiones anteriores descansan en un estudio de simulación con una réplica por escenario, de modo que sus cifras son descriptivas y no estiman la variabilidad entre realizaciones del generador, y el escenario con cambio de régimen está calibrado para favorecer la separación entre regímenes (Sección~\ref{03_07_02_escenarios}), por lo que la magnitud de la ventaja vale para esa configuración y no para cualquier grado de separación. En el escenario con multimodalidad en los scores, el rasgo queda en la ley conjunta y no llega a la marginal de la curva, de manera que el estudio no observa el caso en que la curva siguiente es multimodal por construcción fuera del cambio de régimen.\\

Los comportamientos reparten masa con más generosidad que la ley verdadera, y lo hacen donde el proceso cambia de régimen. En el escenario con cambio de régimen el modelo declara más de un comportamiento en el $30\%$ de los orígenes de prueba, mientras que la ley condicional es multimodal en cerca de un $9\%$ (Sección~\ref{03_08_03_comportamientos}). La diferencia se concentra en los orígenes cuyo estado se acerca al umbral, porque ahí los dos regímenes son posibles y la transición probit de los pesos, más suave que la del generador, reparte la masa entre ambos antes y después de lo que el proceso lo hace, de modo que al modelo le cuesta más decidir cuanto más cerca está el cambio. Esa masa adicional no es un régimen falso sino una opción de baja probabilidad, y su costo es pequeño, ya que las bandas y la cota no la pagan, pero sí afecta al pronóstico por comportamiento como regla de elección. Cuando la probabilidad del comportamiento elegido queda entre $0.5$ y $0.7$, elegir pierde más de lo que gana, y con los pocos orígenes de ese tramo no puede separarse si se debe a estados sobre el umbral, donde el propio proceso es incierto, o a una probabilidad optimista del modelo. Por eso la recomendación práctica es reportar las opciones y sus probabilidades en vez de elegir cuando la probabilidad es baja.\\

En la aplicación, el muestreador se ejecutó con una sola cadena por especificación, de modo que no se calcula el diagnóstico de convergencia entre cadenas y las probabilidades posteriores de inclusión se leen como descriptivas. Las mezclas de los scores menores ocupan casi todo el truncamiento de 30 átomos, y en ellas la probabilidad de inclusión no distingue entre covariables (Sección~\ref{04_03_05_seleccion}). Para que el modelo sea competitivo es necesario retener prácticamente tantos scores como funciones tiene la base, aquí $M = J = 10$, de modo que la reducción de dimensión que motiva el uso de los scores es, en la práctica, pequeña.\\

Por último, la factorización por componentes, que modela cada score con una mezcla univariada y supone independencia condicional entre ellos, es la que hace tratable la inferencia y permite ejecutar los $M$ muestreadores en paralelo, pero condiciona la interpretación. Las autofunciones se extienden sobre todo el dominio, de modo que un score no corresponde a un intervalo particular de la curva y la selección de variables indica qué scores rezagados afectan a cada score sin traducirse en qué horas del día anterior afectan a qué horas del siguiente. Una base local conservaría esa lectura, pero sus coeficientes están correlacionados y no admiten la factorización sin contradecir su supuesto. La interpretabilidad de la selección y la factorización por componentes quedan en tensión, como se discute en la Sección~\ref{04_04_00_discusion}.

\section{Trabajo Futuro}
\label{05_03_00_trabajo_futuro}

Las limitaciones anteriores señalan las direcciones en que el trabajo puede continuar. La primera es completar el estudio de simulación con un escenario en que la curva siguiente sea multimodal por construcción, con regímenes que difieran en nivel y no solo en dispersión, de modo que los comportamientos puedan contrastarse con modos verdaderos de la curva y no solo con los de la ley conjunta de los scores, y con varias réplicas por escenario, para que las diferencias entre modelos lleven una medida de su variabilidad.\\

La segunda es el modelo conjunto sobre el vector de scores. La factorización es un supuesto de conveniencia, y una mezcla única para $\boldsymbol{\xi}_t$ con átomos multivariados, cuyos pesos dependieran del mismo estado rezagado, permitiría que un mismo átomo recogiera la configuración completa de un régimen en vez de repartirla entre $M$ mezclas que solo comparten $\mathbf{z}_t$. Esa formulación eliminaría la tensión entre interpretabilidad y factorización, porque admitiría bases locales con coeficientes correlacionados, al costo de un muestreador de mayor dimensión que ya no se paraleliza por componente y de una identificación de los comportamientos que podría leerse directamente de los átomos. Una alternativa más directa, y previa a esa formulación, es ajustar el modelo tal como está sobre los coeficientes de una base local, aceptando que se viola la independencia condicional entre componentes, para ver cuánto cuesta esa violación en el desempeño predictivo. El interés está en que los scores de la FPCA, y en general los coeficientes de cualquier base ortonormalizada, no heredan la localidad de la curva, mientras que los coeficientes de una B-spline sí la tienen, cada uno gobierna un tramo del dominio; pero esa localidad se pierde en cuanto se corrige por la matriz de Gram para trabajar en la métrica de $L^2$, porque la transformación vuelve a repartir cada coeficiente sobre todo el dominio. Ajustar sobre los coeficientes crudos de la B-spline conservaría la lectura por tramos de la selección de variables, qué horas del día anterior afectan a qué horas del siguiente, y la comparación con el ajuste sobre scores diría si la independencia condicional es un supuesto que el modelo necesita o solo uno que lo hace cómodo.\\

La tercera es la calibración de las probabilidades cerca de un cambio de régimen. La transición probit de los pesos es más suave que la del proceso, de modo que al acercarse el estado al umbral la mezcla reparte la masa entre los regímenes antes y después de lo que corresponde, y eso es lo que separa al pronóstico de su cota en los orígenes de probabilidad intermedia. Una especificación de los pesos que admita transiciones más abruptas, sea con un prior menos restrictivo sobre los coeficientes de decaimiento o con un enlace distinto del probit sobre distancias, acercaría las probabilidades declaradas a las del proceso sin quitar opciones reales. Medirlo requiere más orígenes cerca del umbral que los que da una réplica, para separar la incertidumbre propia del proceso en ese tramo de la que agrega el modelo.\\

La cuarta es ampliar lo que entra al estado. La factorización admite covariables exógenas en $\mathbf{z}_t$ sin alterar su estructura, de modo que el calendario, que en el tráfico se usó solo para leer los resultados, podría entrar a los pesos de la mezcla y resolver los días que ningún rezago anuncia, y lo mismo vale para variables meteorológicas o de mercado en las otras series. Del mismo modo, el estudio se limitó a predicciones a un paso, y la extensión a horizontes mayores requiere propagar la predictiva a través de los rezagos, con comportamientos que se ramifican en cada paso.\\

Por último, el costo computacional. Con rezagos propios basta entrenar el mayor $M$ del barrido y reutilizar sus trazas para los menores, pero con los rezagos de todas las componentes cada $M$ exige un muestreador propio, y en las series reales las mezclas de los scores menores ocupan todo el truncamiento. Un truncamiento adaptativo, varias cadenas en paralelo y diagnósticos sobre cantidades invariantes a la permutación de los átomos harían más seguro el uso del modelo en series largas, que son las que más se benefician de una predictiva que cambia con el estado.
